% Лекц 13. Тодорхойгүй интеграл ба эх функц
% Эх: latex/lectures/lec13-antiderivative.tex (scripts/latex_to_beamer.py-ээр хөрвүүлсэн)
% Frame бүр слайдад нэг (эсвэл allowframebreaks-аар хэд хэдэн) хуудас болно;
% хэвлэх хувилбарт frame-ийн хил харагдахгүй, текст үргэлжилсэн байна.

\MathLecture{13}{Тодорхойгүй интеграл ба эх функц}

\section{Уламжлалын урвуу үйлдэл. Эх функц}
\begin{frame}[allowframebreaks]{Уламжлалын урвуу үйлдэл. Эх функц}
\begin{definition}
 $f(x)$ функц өгөгдсөн байг. Хэрэв $\Phi'(x)=f(x)$ бол $\Phi(x)$-ийг $f(x)$ функцийн
\emph{эх функц} гэнэ.   
\end{definition}
\end{frame}

\begin{frame}[allowframebreaks]{Уламжлалын урвуу үйлдэл. Эх функц}
\begin{example}
    $(x^2+x-1)'=2x+1$ тул $x^2+x-1$ нь $2x+1$ функцийн эх функц юм.
\end{example}
\end{frame}

\begin{frame}[allowframebreaks]{Уламжлалын урвуу үйлдэл. Эх функц}
\begin{theorem}\label{togtmol}
   Хэрэв $\forall x\in I$ дээр $f'(x)=g'(x)$ бол $\forall x\in I$ хувьд $f(x)-g(x)=C$ (тогтмол) байна.
 \end{theorem} 
\begin{prof} $h(x)=f(x)-g(x)$ гэе. 
$f'(x)=g'(x) \Rightarrow h'(x)=(f(x)-g(x))'=0$ болно. Лагранжийн теоремийн мөрдөлгөөгөөр   
 $$h'(x)=0 \Rightarrow h(x)=C.$$ Тэгэхээр
\[
f(x)=g(x)+C.
\]
\end{prof}
\end{frame}

\begin{frame}[allowframebreaks]{Уламжлалын урвуу үйлдэл. Эх функц}
\begin{definition}
 $f(x)$ функц өгөгдсөн байг. Түүний бүх эх функцүүдийн олонлогийг  $f(x)$ функцийн тодорхой биш интеграл гээд 
 $\displaystyle\int f(x)\,dx$ гэж тэмдэглэнэ. Th \ref{togtmol}-с 
 $$\int f(x)\,dx=\Phi(x)+C$$ байна. Энд $\Phi(x)$  ямар нэг эх функц.
 \end{definition}
\end{frame}

\section{Тодорхойгүй интегралын үндсэн жагсаалт}
\begin{frame}[allowframebreaks]{Тодорхойгүй интегралын үндсэн жагсаалт}
\begin{multicols}{2}
\begin{enumerate}
\item $\displaystyle\int x^n\,dx = \frac{x^{n+1}}{n+1}+C,\quad (n\neq -1)$
\item $\displaystyle\int \frac{dx}{x}= \ln |x|+C$
\item $\displaystyle\int e^x dx = e^x +C $
\item $\displaystyle\int a^x dx = \frac{a^x}{\ln a}+C$
\item $\displaystyle\int \sin x dx = -\cos x + C $
\item $\displaystyle\int \cos x dx = \sin x + C$
\item $\displaystyle\int \sec^2 x\, dx = \tan x +C$ 
\item $\displaystyle\int \frac{dx}{1+x^2}= \arctan x + C$
\item $\displaystyle\int \frac{dx}{\sqrt{1-x^2}} = \arcsin x +C$
\end{enumerate}
\end{multicols}
\end{frame}

\begin{frame}[allowframebreaks]{Тодорхойгүй интегралын үндсэн жагсаалт}
\begin{theorem}[ Арифметик үйлдлүүд ба эх функц] $f, \, g$ нь эх функцтэй бол, $C$ тогтмол тоо бол
  \begin{enumerate}
  \item $\displaystyle \int Cf\,dx=C\displaystyle \int f\,dx$ 
\item $\displaystyle \int (f\pm g)\,dx = \int f\,dx \pm \int g\,dx$
\item $\displaystyle \int f'(x)g(x)\,dx = f(x)g(x) - \int f(x)g'(x)\,dx$ (Хэсэгчлэн интегралчлах)
\item $\displaystyle \int f(g(x))g'(x)\,dx = \int f(u)\,du$ Энд $u=g(x)$,\quad (Орлуулах)
  \end{enumerate}
\end{theorem}
\end{frame}

\begin{frame}[allowframebreaks]{Тодорхойгүй интегралын үндсэн жагсаалт}
\begin{example}
\begin{enumerate}
\item $\displaystyle \int 3x^2\,dx=3\displaystyle \int x^2\,dx=3\dfrac{x^{2+1}}{2+1}+c=x^3+c$
\item $\displaystyle \int (3x^2 - 5\sqrt{x^3} + \dfrac{2}{x})\,dx=\displaystyle \int 3x^2 \,dx-\displaystyle \int  5\sqrt{x^3}\,dx+\displaystyle \int\dfrac{2}{x}\,dx=x^3-2x^{2.5}+2\ln x+C$
\item $
\displaystyle\int \sin^2\!\left(\frac{x}{2}\right)\,dx
= \int \frac{1-\cos x}{2}\,dx
= \frac12\int (1-\cos x)\,dx
= \frac{x}{2} - \frac{\sin x}{2} + C.
$
\end{enumerate}  
\end{example}
\end{frame}

\section{Орлуулах арга}
\begin{frame}[allowframebreaks]{Орлуулах арга}
\begin{theorem}[Орлуулах арга]\label{orluulah}
$\displaystyle\int f(u)\,du=F(u)+C$ ба $u=\varphi(x)$ нь дифференциалчлагддаг функц бол
\[
\int f\bigl(\varphi(x)\bigr)\varphi'(x)\,dx=F\bigl(\varphi(x)\bigr)+C .
\]
\end{theorem}
\begin{prof}
Давхар функцийн уламжлалын дүрмээр
$\bigl(F(\varphi(x))\bigr)'=F'(\varphi(x))\,\varphi'(x)=f(\varphi(x))\,\varphi'(x)$.
Иймд $F(\varphi(x))$ нь $f(\varphi(x))\varphi'(x)$ функцийн эх функц болно.
\end{prof}
\end{frame}

\begin{frame}[allowframebreaks]{Орлуулах арга}
$\varphi'(x)\,dx=d\varphi(x)$ тул теоремыг
\[
\int f\bigl(\varphi(x)\bigr)\,d\varphi(x)=\int f(u)\,du\Big|_{u=\varphi(x)}
\]
гэж бичиж болно. Үүнийг \emph{дифференциалын доор оруулах} гэж нэрлэдэг.
Ихэвчлэн хэрэглэгддэг хувиргалтууд:
\[
\renewcommand{\arraystretch}{1.6}
\begin{array}{ll@{\qquad}ll}
dx=\dfrac1a\,d(ax+b)\ (a\ne0), & & \dfrac{dx}{x}=d(\ln|x|), \\
x\,dx=\dfrac12\,d(x^2), & & e^x\,dx=d(e^x), \\
\cos x\,dx=d(\sin x), & & \dfrac{dx}{\cos^2x}=d(\tg x), \\
\sin x\,dx=-d(\cos x), & & \dfrac{dx}{2\sqrt x}=d(\sqrt x).
\end{array}
\]
\end{frame}

\begin{frame}[allowframebreaks]{Орлуулах арга}
\begin{cor}
$\displaystyle\int f(x)\,dx=F(x)+C$ бол
\begin{enumerate}
\item $\displaystyle\int f(ax+b)\,dx=\frac1a\,F(ax+b)+C$, \quad $a\ne0$;
\item $\displaystyle\int\frac{\varphi'(x)}{\varphi(x)}\,dx=\ln|\varphi(x)|+C$.
\end{enumerate}
\end{cor}
\end{frame}

\begin{frame}[allowframebreaks]{Орлуулах арга}
\begin{example}
\begin{enumerate}
\item $\displaystyle\int(3x-1)^5\,dx=\frac13\int(3x-1)^5\,d(3x-1)=\frac{(3x-1)^6}{18}+C$.
\item $\displaystyle\int\cos(2x+1)\,dx=\frac12\sin(2x+1)+C$.
\item $\displaystyle\int\frac{dx}{a^2+x^2}=\frac1a\int\frac{d(x/a)}{1+(x/a)^2}=\frac1a\arctg\frac xa+C$, \ ($a>0$).
\item $\displaystyle\int\frac{dx}{\sqrt{a^2-x^2}}=\int\frac{d(x/a)}{\sqrt{1-(x/a)^2}}=\arcsin\frac xa+C$, \ ($a>0$).
\end{enumerate}
\end{example}
\end{frame}

\begin{frame}[allowframebreaks]{Орлуулах арга}
\begin{example}
\begin{enumerate}
\item $\displaystyle\int x\,e^{x^2}dx$. \ $u=x^2$, $du=2x\,dx$ гэвэл
$\displaystyle\int x\,e^{x^2}dx=\frac12\int e^u\,du=\frac12e^{x^2}+C$.
\item $\displaystyle\int\tg x\,dx=\int\frac{\sin x}{\cos x}\,dx=-\int\frac{d(\cos x)}{\cos x}=-\ln|\cos x|+C$.
\item $\displaystyle\int\sin^3x\cos x\,dx=\int\sin^3x\,d(\sin x)=\frac{\sin^4x}{4}+C$.
\item $\displaystyle\int\frac{\ln x}{x}\,dx=\int\ln x\,d(\ln x)=\frac{\ln^2x}{2}+C$.
\item $\displaystyle\int\frac{2x+3}{x^2+3x+5}\,dx=\int\frac{d(x^2+3x+5)}{x^2+3x+5}=\ln(x^2+3x+5)+C$.
\end{enumerate}
\end{example}
\end{frame}

\begin{frame}[allowframebreaks]{Орлуулах арга}
\begin{theorem}[Хувьсагч солих]\label{huvsagch}
$x=\psi(t)$ нь $T$ завсар дээр дифференциалчлагддаг, хатуу монотон функц бөгөөд $t=\psi^{-1}(x)$ урвуутай байг. Хэрэв
$\displaystyle\int f\bigl(\psi(t)\bigr)\psi'(t)\,dt=G(t)+C$ бол
\[
\int f(x)\,dx=G\bigl(\psi^{-1}(x)\bigr)+C .
\]
\end{theorem}
\begin{prof}
Урвуу функцийн уламжлалаар $\bigl(\psi^{-1}(x)\bigr)'=\dfrac1{\psi'(t)}$. Иймд
$\bigl(G(\psi^{-1}(x))\bigr)'=G'(t)\cdot\dfrac1{\psi'(t)}=f(\psi(t))=f(x)$.
\end{prof}
\begin{remark}
Энэ аргад шинэ хувьсагчийг $x=\psi(t)$ хэлбэрээр сонгоно. Интегралчилсны дараа $t=\psi^{-1}(x)$-ийг орлуулж, анхны хувьсагч руу буцна.
\end{remark}
\end{frame}

\begin{frame}[allowframebreaks]{Орлуулах арга}
\begin{example}
$\displaystyle\int x\sqrt{x-1}\,dx$. \ $t=\sqrt{x-1}$ гэвэл $x=t^2+1$, $dx=2t\,dt$. Иймд
\[
\int x\sqrt{x-1}\,dx=\int(t^2+1)\,t\cdot2t\,dt=\frac{2t^5}{5}+\frac{2t^3}{3}+C
=\frac25(x-1)^{5/2}+\frac23(x-1)^{3/2}+C .
\]
\end{example}
\begin{example}
$\displaystyle\int\frac{dx}{1+\sqrt x}$. \ $x=t^2$ ($t\ge0$), $dx=2t\,dt$ гэвэл
\[
\int\frac{2t\,dt}{1+t}=2\int\Bigl(1-\frac1{1+t}\Bigr)dt=2t-2\ln(1+t)+C=2\sqrt x-2\ln(1+\sqrt x)+C .
\]
\end{example}
\end{frame}

\begin{frame}[allowframebreaks]{Орлуулах арга}
\begin{example}
$\displaystyle\int\sqrt{1-x^2}\,dx$. \ $x=\sin t$, $t\in[-\frac\pi2,\frac\pi2]$ гэвэл $dx=\cos t\,dt$, $\sqrt{1-x^2}=\cos t$. Иймд
\[
\int\sqrt{1-x^2}\,dx=\int\cos^2t\,dt=\int\frac{1+\cos2t}{2}\,dt=\frac t2+\frac{\sin2t}{4}+C .
\]
$\sin2t=2\sin t\cos t=2x\sqrt{1-x^2}$ тул
\[
\int\sqrt{1-x^2}\,dx=\frac12\arcsin x+\frac x2\sqrt{1-x^2}+C .
\]
\end{example}
\begin{remark}
$\sqrt{a^2-x^2}$ агуулсан интегралд $x=a\sin t$, \ $\sqrt{a^2+x^2}$ агуулсан интегралд $x=a\tg t$,
\ $\sqrt[n]{ax+b}$ агуулсан интегралд $t=\sqrt[n]{ax+b}$ орлуулгыг ихэвчлэн хэрэглэнэ.
\end{remark}
\end{frame}

\section{Хэсэгчлэн интегралчлах арга}
\begin{frame}[allowframebreaks]{Хэсэгчлэн интегралчлах арга}
$\displaystyle \int g'(x)f(x)\,dx = f(x)g(x) - \int g(x)f'(x)\,dx$ томьёог
\[
\int fdg = fg - \int g\,df
\] гэж бичиж болно. 

\textbf{Жишээ:}
$\int x e^x dx $ интегралд $f(x)=x, g'(x)dx=dg(x)=e^xdx=de^x$ гэвэл
\[
\int x e^x dx = xe^x - \int e^x dx = e^x(x -1)+C
\] болно. Хэрэв энд $f(x)=e^x, g'(x)dx=dg(x)=xdx=d\frac{x^2}{2}$ гэвэл
\[
\int x e^x dx = \frac{x^2}{2}e^x - \int \frac{x^2}{2}e^x dx
\] болно. Энэ гарсан интеграл нь өмнөхөөс хүнд харагдаж байна. Иймд хэсэгчлэх арга хэрэглэхэд $f, g'$ хэрхэн сонгох нь чухал юм.

\textbf{$f(x)$ сонгох дараалал (LAPTE):} Хэсэгчлэх арга хэрэглэхэд ихэвчлэн $f(x)$ функцийг дараах дараалалаар сонгоно.
\[
\text{Log} \rightarrow \text{Arc} \rightarrow \text{Polynomm} \rightarrow \text{Trig} \rightarrow \text{Exponential}
\]
\end{frame}

\begin{frame}[allowframebreaks]{Хэсэгчлэн интегралчлах арга}
\begin{example}
    $\displaystyle\int x\ln x \,dx$ бодоход LAPTE сонголтоор $\ln x$ байгаа тул түүнийг $f(x)=\ln x, g'(x)=x$ гэж авна. Иймд $g(x)=\dfrac{x^2}{2}$ болно. Иймд 
    $$\int x\ln x \,dx=\dfrac{x^2}{2}\ln x-\int \dfrac{x^2}{2}(\ln x)'dx=\dfrac{x^2}{2}\ln x-\int \dfrac{x^2}{2}\dfrac{1}{x}dx=\dfrac{x^2}{2}\ln x-\dfrac{x^2}{4}+c$$ болно.
\end{example}
\end{frame}

\begin{frame}[allowframebreaks]{Хэсэгчлэн интегралчлах арга}
\begin{example} $\int x^{2}e^{x}\,dx$ бодоход LAPTE сонголтоор $f(x)=x^2, g'(x)=e^x$ гэж авна. Иймд $g(x)=e^x$ болно.
 \[
\int x^{2}e^{x}\,dx
= x^{2}e^{x}-\int 2xe^{x}\,dx.
\] болно.
Сүүлийн $\displaystyle\int 2xe^{x}\,dx$ интегралыг мөн хэсэгчлэн интегралчилбал 
$$\int 2xe^{x}\,dx= 2xe^{x}-\int 2e^{x}\,dx$$ болно. Иймд 
\[
\int x^{2}e^{x}\,dx= x^{2}e^{x}-2xe^{x}+2e^{x}+C
= e^{x}(x^{2}-2x+2)+C.
\]
болно.
\end{example} 
\end{frame}

\begin{frame}[allowframebreaks]{Хэсэгчлэн интегралчлах арга}
\begin{example} $\displaystyle\int e^{x}\sin x\,dx $ интегралын хувьд  $f(x)=\sin x, g'(x)=e^x$ гэж авбал
    \[
\int e^{x}\sin x\,dx =e^x\sin x-\int e^x (-\cos x) dx
\]
болно. $\int e^x (-\cos x) dx$ интегралын хувьд  $f(x)=-\cos x, g'(x)=e^x$ гэж авбал

$$\int e^x (-\cos x) dx=e^x (-\cos x)-\int e^x(-\sin x) dx$$ болно. Хэрэв $I=\displaystyle\int e^{x}\sin x\,dx$ гэвэл 
$I=e^x\sin x-e^x\cos x-I$ болох тул $I=\dfrac{e^{x}}{2}(\sin x - \cos x) + C.$
\end{example}
\end{frame}

% Семинар ба гэрийн даалгавар (Берманы бодлогын хураамжаас)
% Лекц 13: семинар ба гэрийн даалгавар. Бодлогууд: Г.Н. Берман, "Сборник задач по курсу
% математического анализа" (22-р хэвлэл). [Б.N] нь номын бодлогын дугаар.

\ProblemSet{Семинар 13}
\begin{frame}[allowframebreaks]{Семинар 13}
\pgroup{Үндсэн жагсаалт ба энгийн хувиргалтаар интегралыг ол.}
\begin{plist}
\prob{1676}{$\int\sqrt x\,dx$}
\prob{1678}{$\int\frac{dx}{x^2}$}
\prob{1679}{$\int10^x\,dx$}
\prob{1685}{$\int(\sqrt x+1)(x-\sqrt x+1)\,dx$}
\prob{1688}{$\int\bigl(\frac{1-z}{z}\bigr)^2dz$}
\prob{1694}{$\int\frac{1+\cos^2x}{1+\cos2x}\,dx$}
\prob{1696}{$\int\tg^2x\,dx$}
\end{plist}
\pgroup{Дифференциалын доор оруулах (хувьсагч солих) аргаар ол.}
\begin{plist}
\prob{1706}{$\int(x+1)^{15}dx$}
\prob{1710}{$\int\sqrt{8-2x}\,dx$}
\prob{1713}{$\int x\sqrt{1-x^2}\,dx$}
\prob{1719}{$\int\sin^3x\cos x\,dx$}
\prob{1744}{$\int\tg x\,dx$}
\prob{1749}{$\int\frac{dx}{x\ln x}$}
\prob{1760}{$\int\frac{dx}{1+9x^2}$}
\prob{1761}{$\int\frac{dx}{\sqrt{4-x^2}}$}
\end{plist}
\pgroup{Хэсэгчлэн интегралчил.}
\begin{plist}
\prob{1833}{$\int x\cos x\,dx$}
\prob{1834}{$\int xe^{-x}dx$}
\prob{1836}{$\int x^n\ln x\,dx$ \ ($n\ne-1$)}
\prob{1837}{$\int x\arctg x\,dx$}
\prob{1860}{$\int e^x\sin x\,dx$}
\end{plist}
\end{frame}

\ProblemSet{Гэрийн даалгавар 13}
\begin{frame}[allowframebreaks]{Гэрийн даалгавар 13}
\pgroup{Үндсэн жагсаалт ба энгийн хувиргалтаар интегралыг ол.}
\begin{plist}
\prob{1681}{$\int\frac{dx}{2\sqrt x}$}
\prob{1683}{$\int3{,}4\,x^{-0{,}17}dx$}
\prob{1689}{$\int\frac{(1-x)^2}{x\sqrt x}\,dx$}
\prob{1690}{$\int\frac{(1+\sqrt x)^3}{\sqrt[3]{x}}\,dx$}
\prob{1693}{$\int\frac{3\cdot2^x-2\cdot3^x}{2^x}\,dx$}
\prob{1695}{$\int\frac{\cos2x}{\cos^2x\sin^2x}\,dx$}
\prob{1697}{$\int\ctg^2x\,dx$}
\prob{1698}{$\int2\sin^2\frac{x}{2}\,dx$}
\prob{1699}{$\int\frac{(1+2x^2)\,dx}{x^2(1+x^2)}$}
\end{plist}
\pgroup{Дифференциалын доор оруулах (хувьсагч солих) аргаар ол.}
\begin{plist}
\prob{1707}{$\int\frac{dx}{(2x-3)^5}$}
\prob{1712}{$\int2x\sqrt{x^2+1}\,dx$}
\prob{1715}{$\int\frac{x\,dx}{\sqrt{x^2+1}}$}
\prob{1720}{$\int\frac{\sin x\,dx}{\cos^2x}$}
\prob{1723}{$\int\frac{\sqrt{\ln x}}{x}\,dx$}
\prob{1724}{$\int\frac{(\arctg x)^2dx}{1+x^2}$}
\prob{1731}{$\int\sin(2x-3)\,dx$}
\prob{1738}{$\int\frac{dx}{2x-1}$}
\prob{1742}{$\int\frac{e^x\,dx}{e^x+1}$}
\prob{1752}{$\int e^{\sin x}\cos x\,dx$}
\prob{1755}{$\int e^{-3x+1}dx$}
\prob{1763}{$\int\frac{dx}{\sqrt{4-9x^2}}$}
\prob{1774}{$\int\frac{3x-1}{x^2+9}\,dx$}
\end{plist}
\pgroup{Хэсэгчлэн интегралчил.}
\begin{plist}
\prob{1832}{$\int x\sin2x\,dx$}
\prob{1838}{$\int\arccos x\,dx$}
\prob{1842}{$\int x\cos^2x\,dx$}
\prob{1849}{$\int x^2\ln(1+x)\,dx$}
\prob{1850}{$\int x^2e^{-x}dx$}
\prob{1855}{$\int\ln^2x\,dx$}
\prob{1858}{$\int(\arcsin x)^2dx$}
\prob{1862}{$\int e^{ax}\cos nx\,dx$}
\end{plist}
\end{frame}

